\section{\Veche{}}
\label{sec:model}

Consider a set $[E] \triangleq \set{1, 2, \ldots, E}$ of $E$ entities 
that may be recommended to Alice.
We consider the problem of creating a feed $f \in \cF^E$,
where
\begin{align}
    \cF^E \triangleq \set{
        f \in [E]^S \st 
        S \in \setN \wedge \card{f([S])} = S
    }.
\end{align}
Put differently, a feed $f$ is a ranked subset of $[E]$.
Denoting $\Image(f) \subset [E]$ the recommended subset,
$f$ must define an order of $\Image(f)$, 
i.e. a bijection $\card{\Image(f)} \to \Image(f)$.
Let $[U]$ be the set of the $U$ users that Alice follows,
which we call the \emph{councelors}.
As stated in the introduction, 
we propose to construct the feed $f$
by making the $U$ councelors vote 
for what entities they want to recommend to Alice.
More precisely, a stochastic feed based on councelors' votes must be a map
\begin{align}
    \textsc{Ballot}^U \to \Delta(\cF^E),
\end{align}
where $\textsc{Ballot}$ is a set of allowed ballots
and where $\Delta(\cF^E)$ is the set of distributions over $\cF^E$.

\Veche{} is a stochastic feed based on councelors' votes,
where the set of ballots is $\textsc{Ballots} \triangleq \setR^E$.
Put differently, in \Veche{}, each councelor $u$ may submit any ballot $b_u \in \setR^E$,
where $b_{ue} \in \setR$ corresponds to the extent to which 
the councelor $u$ wants entity $e \in [E]$ to be included into Alice's feed.
Furthermore, \Veche{} uses an intermediary ballot aggregation
that maps the councelors' ballots $b \in \setR^{U \times E}$
to a weight vector $w \in \setR^E_{\geq 0}$.
We then leverage feed algorithms that guarantee
that the probability of recommending $e \in [E]$ is tightly linked to the weight $w_e$.
In principle, \Veche{} could include any ballot vector aggregation rule
$\setR^{U \times E} \to \setR^E_{\geq 0}$ to transform councelors' ballots into an entity weight vector.
It is important, however, that the mechanism be mathematically well-behaved.
Typically, it should bound the influence of any councelor $u \in [U]$.


\subsection{\SimpleVeche{}: \Veche{} in its simplest form}

For the sake of exposition,
let use first state a simple instantiation of \Veche{},
which we call \SimpleVeche{}.
We will complexify it in Section~\ref{sec:vote_generalization}.
\SimpleVeche{} first considers the simple ballots given by
\begin{align*}
    b_{ue} \triangleq \left\{
    \begin{array}{ll}
        1, & \text{if } u \text{ posted or reposted } e, \\
        -1, & \text{if } u \text{ reported } e, \\
        0, & \text{otherwise.}
    \end{array}
    \right.
\end{align*}
Then, \SimpleVeche{} simply normalize s
each councelor's ballot $b_u$ by the some of its weights, i.e.
\begin{align}
    \tilde b_u \triangleq \left\{
    \begin{array}{ll}
        b_u / \norm{b_u}{1}, & \text{if } \norm{b_u}{1} > 0, \\
        0, & \text{otherwise}.
    \end{array}
    \right.
\end{align}
Next, \SimpleVeche{} may simply add up all councelors' normalized ballots,
while demanding nonnegative weights, yielding
\begin{align}
    \Weights(b) \triangleq \max \set{0, \sum_{u \in [U]} \tilde{b}_u }.
\end{align}
Finally, \SimpleVeche{} constructs the feed $f$
by repeated sampling without replacement,
where the probability of selecting $e$ is proportional to $w_e$ at each sampling.
\SimpleVeche{} is detailed in Algorithm~\ref{alg:base_veche}.

\begin{algorithm}
\caption{\SimpleVeche{}: A simple instantiation of \Veche{}}\label{alg:base_veche}
\begin{algorithmic}
\Require Maximal size $S$ of returned recommendation subsets. Councelors' ballots $b \in \setR^{U \times E}$.
\State $\tilde b_{u} \gets b_u / \norm{b_u}{1}$
    \Comment{Normalization (with convention $0 / 0 = 0$)}
\State $w \gets \max \set{0, \sum_{u \in [U]} \tilde b_{u} }$
    \Comment{Aggregation of councelors' votes}
\State $\bE \gets \set{e \in [E] \st w_e > 0}$
    \Comment{Filters non-recommendable entities}
\State $f \gets (\perp, \ldots, \perp) \in \set{\perp}^\bE$
    \Comment{Instantiate recommendation subset}
\For{$i = 1, 2, \ldots, \card{\bE}$}
    \Comment{Subset sampling without replacement}
    \State $p \gets w / \norm{w}{1}$
        \Comment{Normalize weights to obtain a probability}
    \State $e^* \gets \Choice(p)$
        \Comment{Select entity $e$ with probability $p_e$}
    \State $f_i \gets e^*$
        \Comment{Append selected entity to ranking}
    \State $w_{e^*} \gets 0$
        \Comment{No replacement: prevent reselecting selected entity}
\EndFor \\
\Return Feed $f$ of the subset $\bE$ of recommendable entities.
\end{algorithmic}
\end{algorithm}

Overall, \SimpleVeche{} defines a map $\setR^{U \times E} \to \Delta (\cF^E)$,
which inputs councelors' ballots $b \in \setR^{U \times E}$
and returns a probability distribution over rankings 
of a subset of recommendable entities in $[E]$.


\subsection{Main gurantee}

Let us formalize the main guarantees of \SimpleVeche{},
which essentially consist of bounding the maximal impact of a councelor
on Alice's feed.
As Alice's feed is random, we first need to a metric over probability distribution.
In this paper, we will use the classical total variation distance.

\begin{definition}
    The total variation distance between
    the probability distributions $\tilde x, \tilde y \in \Delta(\bbX)$
    over $\bbX$ is the maximal probability distance 
    over measurable subsets $X \subset \bbX$, i.e.
    \begin{align*}
        \TotalVariationDistance (\tilde x, \tilde y)
        \triangleq \sup_{X \subset \bbX \text{ measurable}} \absv{
            \pb{x \in X \st x \sim \tilde x}
            - \pb{y \in X \st y \sim \tilde y}
        }.
    \end{align*}
\end{definition}

To state our guarantee, 
we also consider the following indicator,
which we show to be a seminorm.

\begin{definition}
Define the minimal remaining weights at the $k$-th recommendations as
\begin{align}
    \norm{w}{1,k}
    \triangleq \inf_{\substack{\bE \subset [E] \\ \card{\bE} = E - k + 1}} 
    \sum_{e \in \bE} \absv{w_e}.
\end{align}
Note that the minimal weight after $0$ recommendations matches the $\ell_1$ norm,
i.e. $\norm{w}{1,1} = \norm{w}{1}$.
\end{definition}

\begin{proposition}
    \label{prop:seminorm}
    $\norm{\cdot}{1,k}$ is a seminorm.
\end{proposition}

\begin{proof}
    The proof is given in Appendix~\ref{app:seminorm}.
\end{proof}

Denote $\SimpleVeche_{1:k} (b) \in \Delta\left( \left([E] \cup \set{\perp}\right)^k \right)$ 
the distribution 
of the subset of entities in the top $k$ recommendations.
In particular, $e \in \SimpleVeche_{1:k} (b)$ is the binary event that says whether $e$
was included in the recommendations.
For $b \in \setR^{U \times E}$ and $u \in [U]$,
denote $b^{-u}$ the ballot obtained by removing councelor $u \in [U]$.
Then the following holds.

\begin{theorem}
\label{th:simpleveche}
    Let $b \in \setR^{U \times E}$ and $k \in \setN$.
    If $\norm{W(b)}{1, k} > 0$, then
    \begin{align}
        \TotalVariationDistance \left(
            \SimpleVeche_{1:k} (b),
            \SimpleVeche_{1:k} (b^{-u})
        \right)
        \leq k \cdot \norm{W(b)}{1, k}^{-1}.
    \end{align}
\end{theorem}

\begin{proof}
    The proof is given in Appendix~\ref{app:simpleveche}.
\end{proof}

The resilience guarantee may feel underwhelming, 
as it depends on the cumulative nonnegative weights assigned in the absence of councelor $u$.
In fact, this highlights the fact that the vulnerability of the recommendation system is increased
if most councelors spend their ballots to oppose the recommendation of some entities,
thereby leaving out too few votes in favor of entity recommendation.
To clarify the role of negative ballot weights on the security guarantees
with respect to a malicious councelor,
we derive a corollary of the theorem for nonnegative ballot weights.

\begin{corollary}
    Assume nonnegative ballots $b \in \setR^{U \times E}_{\geq 0}$.
    Then
    \begin{align}
        \TotalVariationDistance \left( 
            \SimpleVeche_1(b), 
            \SimpleVeche_1(b^{-u})
        \right) 
        \leq 1 / \left( U - 1 \right).
    \end{align}
\end{corollary}

\begin{proof}
    Assuming nonnegative votes, we have 
    $\norm{W(b^{-u})}{1} 
    = \sum_{e \in [E]} \sum_{v \neq u} \tilde b_{ve} 
    = \sum_{v \neq u} \norm{b_v}{1} = \sum_{v \neq u} 1 
    = U - 1$.
    We then apply Theorem~\ref{th:simpleveche} with $k=1$.
\end{proof}

% \len{
%     I don't yet quite understand the implications of the pivot scheduler, 
%     but intuitively, I feel that weird stuffs can happen 
%     because of the joint probabilities
%     of two entities $e$ and $f$ to be included in $\bS$.
%     I would definitely be interested in better understanding what's happening.
% }

% Denote $\Pivot_e(b, \Scheduler(b))$ the distribution of the event $e \in \bS$
% when running the pivot algorithm on entry $b$ with $\Scheduler$.

% \begin{conjecture}
%     For any ballots $b \in \setR^{U \times E}$,
%     then the total variation distance of $\Pivot(b)$,
%     \begin{align}
%         \sum_{e \in [E]} \TotalVariationDistance \left(
%             \Pivot_e(b, \Scheduler(b)), 
%             \Pivot_e(b_{-u}, \Scheduler(b_{-u})) 
%         \right) \leq 1.
%     \end{align}
% \end{conjecture}


